\scanpage{105}
\mspage{20}
% archive sheet number: 104
% source lines 1--6
\uline{Corollaire} 4.3 Soit $X$ un \deleted{\unsure{sous-}}schéma \deleted{fermé} régulier \deleted{\unsure{au th. de pureté}} \inserted{satisfaisant à la} pureté coh. (4.1), $Y$ un sous-schéma fermé régulier purement de codim. $d$, $n$ un entier premier aux caractéristiques résiduelles de $X$, $\Lambda=\unsure{\mathbb Z/n\mathbb Z|Z/nZ}$, $F$ un faisceau loc. constant sur $X$ annulé par $n$.
\inserted{Pour tout faisceau de $\Lambda$-modules $G$ sur $X$ et tt $d\in\mathbb Z$, on pose $G(d)=G\otimes\mu_n^{\otimes(d)}$.}
\deleted{Alors} on a une suite exacte canonique

\[
\cdots\unsure{{-}|\longrightarrow}H^{i-2d}(Y,F\deleted{\unsure{\otimes\mu_n^{\otimes(-d)}}}(-d))
\xrightarrow{\alpha}H^i(X,F)\xrightarrow{\beta}H^i(X-Y,F)
\xrightarrow{\gamma}H^{i-2d+1}(Y,F\deleted{\unsure{\otimes\mu_n^{\otimes(-d)}}}(-d))\longrightarrow\cdots
\tag*{(4.3.1)}
\]

% source lines 8--9
où $\beta$ est l'homomorphisme de restriction, et où l'homomorphisme composé

\[
H^{i-2}(X,F\deleted{\unsure{\otimes\mu_n^{\otimes d}}}(d))
\xrightarrow{\text{restr.}}H^i(Y,F\deleted{\unsure{\otimes\mu_n^{\otimes(-d)}}}(\unsure{-d|d}))
\xrightarrow{\alpha}H^i(X,F)
\]

% source lines 11--12
est le cup-produit avec une classe de cohomologie \deleted{bien déterminée} \inserted{canonique}

\[
\gamma(Y)\in\deleted{\unsure{\text{de}}}H^{2d}(X,\deleted{\unsure{\mu_X}}\Lambda_X(d)).
\tag*{(4.3.2)}
\]

% source lines 14--17 and additions
\deleted{Quand $Y$ est \unsure{de}} Lorsque $d=1$, \inserted{\unsure{de sorte que} $Y$ est un \uline{diviseur} sur $X$,} cette classe est aussi la \inserted{première} classe de Chern de $\underline L=\unsure{\mathscr O_X|\mathcal O_X}(Y)$, i.e. l'image \inserted{de la classe de $\underline L$ dans $\operatorname{Pic}(X)\simeq H^1(X,\unsure{\mathbb G_{mX}|G_{mX}})$} par l'homomorphisme \deleted{\unsure{de}} de cobord associé à la suite exacte de Kummer

\[
0\longrightarrow\mu_{nX}\longrightarrow\mathbb G_{mX}\overset{\lambda\mapsto\lambda^n}{\unsure{{-}|\longrightarrow}}\mathbb G_{mX}\longrightarrow 0.
\]

% source lines 19--20
\uline{Dém.} La suite exacte (4.3.1) n'est autre essentiellement que la suite exacte de cohomologie relative

\[
\cdots\longrightarrow H^i_Y(X,F)\unsure{{-}|\longrightarrow}H^i(X,F)\unsure{{-}|\longrightarrow}H^i(X-Y,F)\longrightarrow\cdots,
\]

% source line 22
compte tenu des isomorphismes

\[
H^i_Y(X,F)\simeq H^{i-2d}(Y,F\deleted{\unsure{\otimes\mu_n^{\otimes(-d)}}}(-d))
\]

% source line 24
déduits de la suite spectrale de localisation

\[
H^{\unsure{*|n}}_Y(X,F)\Longleftarrow\unsure{E|F}^{pq}_2=H^p(Y,\underline H^q_Y(F)),
\]
